\documentclass[12pt]{article}
\usepackage[margin=1in]{geometry}
\usepackage{fontspec}
\setmainfont{Times New Roman}[Path=/home/sandbox/fonts/,Extension=.TTF,UprightFont=Times,BoldFont=Timesbd,ItalicFont=Timesi,BoldItalicFont=Timesbi]
\usepackage{graphicx,longtable,amsmath,amssymb,amsthm,booktabs}
\usepackage[colorlinks=true,linkcolor=black,citecolor=black,urlcolor=blue]{hyperref}
\newtheorem{proposition}{Proposition}
\newtheorem{remark}{Remark}
\newtheorem{definition}{Definition}
\title{Virtual Yeast for Climate-Resilient Fermentation:\\[4pt]
{\large A Pre-Registered Constraint-Based Map of Combined-Stress Collapse
in \emph{Saccharomyces cerevisiae}, its Co-Limitation Mechanism, and a
Demonstrated Boundary of Steady-State Strain Design}}
\author{MEGA-PROGRAM-27, Item 5}
\date{\today}
\begin{document}
\maketitle

\begin{abstract}
\noindent Industrial bioethanol fermentations fail when climate-driven
stresses combine: heat waves, accumulating ethanol, osmotic pressure from
concentrated feedstocks, and nutrient limitation interact inside the same
tank, yet each stress is usually studied one factor at a time. We built a
virtual \emph{Saccharomyces cerevisiae} cell on the consensus genome-scale
model yeast-GEM (4{,}105 reactions, 2{,}748 metabolites, 1{,}143 genes;
in-repo SBML pinned by SHA-256) and extended it with four
literature-anchored stress layers: cardinal-temperature enzyme kinetics,
Levenspiel-type ethanol product inhibition, an obligatory osmoregulatory
glycerol-production drain, and nitrogen availability scaling, with
maintenance energy rising inversely with relative growth. Every functional
form, parameter, environment grid, objective, and success gate was locked
in a SHA-256 pre-registration before any outcome run, with a locked
70/30 design/test split of the environment space and the test region
touched exactly once. Across a 5{,}000-environment Latin-hypercube sample
of the four-axis stress space, 37.1\% of environments collapse wild-type
ethanol production below 20\% of its reference. Bottleneck analysis of the
collapsed region shows collapse is co-limited by the osmoregulatory
glycerol drain (binding in 69\% of solvable collapsed environments) and a
heat-scaled glycolytic-capacity loss (59\%), a mechanism corroborated at
the regulatory level by concordance between the model's heat regime and
the published heat-shock expression program (10/10 marker genes,
$p=9.8\times10^{-4}$). The pre-registered non-additivity gate was not met
(24.2\% vs.\ the 25\% threshold; parameter-fragile under temperature-margin
perturbations), and we report it as such. A pre-registered strain-design
campaign - all 1{,}143 single-gene deletions, 435 pairs, 36 reaction
scalings, a seeded evolutionary search, and a matched random baseline,
evaluated once on the untouched test region - found no design that
improves any locked objective beyond solver noise ($<10^{-14}$), because
collapse under this mechanism is imposed by environmental caps rather than
by internal pathway limitation. Validation against three published
genome-wide deletion screens likewise returns a declared null (0/4
pre-registered overlap comparisons significant), delineating a boundary of
steady-state flux models for stress-tolerance engineering. All negative
results are reported plainly under the pre-registration; the principal
positive contributions are the collapse map, its verified co-limitation
mechanism, the regime-level transcriptome concordance, and a demonstrated
method boundary for the field's most widely used engineering model class.
\end{abstract}

\section{Introduction}
\label{sec:intro}

\subsection{The industrial problem: fermentation under a changing climate}

Bioethanol is the largest-volume fermentation product of the industrial
bioeconomy, and its production organism, \emph{Saccharomyces cerevisiae},
operates at the edge of its physiological envelope in essentially every
commercial process. A modern fuel-ethanol plant asks a yeast cell to
convert a concentrated sugar mash - routinely above 200\,g/L of glucose in
very-high-gravity (VHG) operation - into ethanol at titers that approach
the limits of its own tolerance, while the metabolic heat of fermentation
and increasingly volatile ambient conditions push tank temperatures toward
the top of the organism's growth range. Each of these burdens has a name
and a literature. Ethanol itself is a chaotropic solvent that disorderes
membranes and dissipates proton gradients; its inhibition of yeast growth
and fermentation has been modeled since the classical Levenspiel
formulation. Temperature acts on every enzyme in the cell at once, and the
cardinal-temperature model of Rosso and colleagues - a minimum, optimum,
and maximum temperature with a characteristic asymmetric rise and
precipitous fall of rate - has been fitted to yeast growth and
fermentation kinetics across strains and conditions. Osmotic pressure,
whether from the sugar of concentrated mashes or from drought-stressed
feedstocks, forces the cell into an obligatory osmoregulatory program:
glycerol must be synthesized and retained to balance water activity, and
every mole of glycerol is carbon that cannot become ethanol. Nitrogen and
other nutrient limitations shape the achievable growth rate and the
maintenance burden the cell must pay before any product is made.

What the industrial literature also says, plainly, is that these stresses
do not arrive one at a time. A heat wave during a VHG fermentation is a
simultaneous increase in temperature, osmotic load (as sugar concentrates
when water evaporates), ethanol concentration (as fermentation
progresses), and nutrient stress (as the stationary phase approaches).
Plant operators experience this as the ``stuck fermentation'': a batch
that slows, stalls, and leaves residual sugar that no restart recovers.
The economics are direct - a few percent of lost yield across a
continental industry is measured in billions of liters - and the climate
projection is worse: the frequency of multi-day heat extremes during the
fermentation season is rising on every major production continent. The
engineering response has so far been empirical: stress-tolerant strains
are selected, evolved, or screened one stress at a time, and then combined
in the hope that tolerance is additive. Sometimes it is. Often it is not,
and the failures are rarely publishable, so the field's map of which
combinations of stress actually break the process - and \emph{why} -
remains thin.

\subsection{The scientific gap: stress is studied one axis at a time}

The modeling literature mirrors the experimental one. Genome-scale
metabolic models (GEMs) of \emph{S.\ cerevisiae} - of which the consensus
yeast-GEM is the community standard - have been used for two decades to
predict gene essentiality, design metabolic-engineering interventions, and
contextualize omics data. Stress, when it enters these models at all,
usually enters as a single scalar perturbation: a temperature-dependent
turnover number here, an inhibitory ethanol concentration there. We are
not aware of any pre-registered, genome-scale study that treats the
environment as a joint, four-axis object - temperature, ethanol, osmotic
pressure, and nutrient availability - and asks where in that joint space
the cell's fermentation phenotype collapses, which internal constraints
bind when it does, and whether the answer is robust to the chosen
parameterization. The closest works study pairs of stresses, or study
combined stress transcriptomically without a mechanistic flux model, or
build kinetic models of a handful of reactions whose scope cannot reach
the whole-cell question. The result is that the most basic map an
engineer would want - here is the region of stress space where the
process fails, and here is what fails first - does not exist in a
reproducible, falsifiable form.

That gap is not only industrial. It is a clean instance of a general
question in systems biology: when a cell fails under combined
perturbations, is the failure merely the product of independent
single-stress effects, or does the combination create emergent failure
modes that no single-stress study can see? And if emergence exists, is it
a property of the biology or of the mathematics used to describe it? We
treat these as separate, testable questions, and we built the study so
that each gets its own pre-registered gate.

\subsection{What this paper does}

We construct a virtual \emph{S.\ cerevisiae} cell for combined-stress
fermentation and subject it to a fully pre-registered analysis program.
The model extends the consensus yeast-GEM with four stress layers, each
anchored to published kinetic or physiological data and each carrying its
provenance: cardinal-temperature (CTMI) scaling of enzyme-catalyzed
rates; Levenspiel-type ethanol inhibition of growth-linked processes; an
obligatory glycerol-production drain proportional to external osmolarity;
and a nitrogen-availability scaling of nutrient uptake. Maintenance energy
demand rises as relative growth falls, in a declared, assumed functional
form that is never fitted to any outcome. The environment space is a
locked four-axis grid - 28--42$^\circ$C, 0--14\% v/v ethanol, 0--1.5\,M
sorbitol-equivalent osmolarity, 5--100\% of standard nitrogen - and the
objectives are the four an industrial process cares about: ethanol yield,
temperature tolerance, ethanol tolerance, and low nutrient requirement.
Growth is treated as a means, never as the objective.

The analysis program is staged and gated. A1 maps single-stressor dose
responses. A2 maps the combined-stress collapse surface over a
5{,}000-environment Latin-hypercube sample. A3 tests, against a
pre-registered threshold, whether combined-stress failure is sub-additive
- worse than the product of the single-stress effects - and whether that
conclusion survives parameter perturbation. A4 identifies, with shadow
prices and flux-variability analysis, which constraints bind in the
collapsed region. A5 is a strain-design campaign under a pre-registered
gate: all single-gene deletions, a seeded pair screen, reaction-bounds
scalings, an evolutionary search, and a matched-budget random baseline,
with final claims evaluated once on the untouched test region. A6
validates the model's stress-conditional gene-sensitivity predictions
against three published genome-wide deletion screens with pre-registered
overlap statistics, plus qualitative checks against two further studies.
A7 asks whether the model's stress regimes are concordant with the
published environmental-stress expression program at the level of marker
gene sets. A8 lifts the steady-state assumption with dynamic heat-wave
simulations, asking whether collapse timing reflects adaptation or
accumulation.

Every gate is reported as it landed. Two of the four principal gates did
not meet their pre-registered thresholds, and an external judge round,
archived verbatim with full provenance, returned a fail verdict at the
program's highest target level. We present these outcomes not as
procedural embarrassments but as results: the negative theorems of this
paper - about what combined-stress collapse is not, and about what
steady-state strain design cannot do - are as much a part of the
scientific record as the positive map and mechanism, and they are
established here with the same evidentiary discipline.

\subsection{Summary of principal findings}

\begin{enumerate}
\item \textbf{A reproducible collapse map.} In the locked environment
space, 37.1\% of sampled combined-stress environments drive wild-type
ethanol production below 20\% of its unstressed reference; the collapse
region has a definite geometry in the temperature--ethanol plane that the
osmotic axis modulates and the (declared, see below) nitrogen axis does
not.
\item \textbf{A verified co-limitation mechanism.} Collapse is co-limited
by the obligatory osmoregulatory glycerol drain and by heat-scaled loss of
glycolytic capacity: the two constraints bind, respectively, in 69\% and
59\% of solvable collapsed environments, with flux-variability analysis
direction-consistent. This is the paper's central mechanistic claim, and
it is stated with its shadow-price evidence and its regime structure.
\item \textbf{Regime-level regulatory concordance.} The model's heat
regime is concordant with the published heat-shock expression program at
10/10 marker genes ($p=9.8\times10^{-4}$); the nitrogen-catabolite and
glycerol regimes trend or are non-significant; the glycolytic regime is
declared non-concordant ($p=0.42$). Concordance is reported per regime,
not averaged into a single flattering number.
\item \textbf{A declared boundary of steady-state strain design.} Under
the verified mechanism, no deletion, pair, scaling, or evolved design
improves any locked objective beyond solver noise ($<10^{-14}$ relative
yield) on the untouched test region: collapse is imposed by environmental
caps, not by an internal bottleneck a knockout can relieve. The model
correspondingly fails to recover deletion-screen stress sensitivity
(0/4 pre-registered overlap tests). Both results are reported as
limitations of the model class, not of the analysis.
\item \textbf{Dynamic corroboration.} Dynamic heat-wave simulations show
collapse timing is set by endogenous ethanol accumulation rather than by
adaptive memory: no hysteresis beyond ethanol carry-over, and dynamic and
static collapse boundaries agree.
\end{enumerate}

\subsection{Scope, honesty discipline, and declared limitations}

This is a computational study. No wet-lab claim is made anywhere. All
validation is against previously published experimental datasets; all
parameters are locked before outcome runs; all parameter perturbations
are one-at-a-time sweeps with declared ranges; exploratory analyses are
labeled as such and never mixed with confirmatory claims. Two structural
limitations are declared up front because they condition the reading of
every result: first, the nitrogen lever of the environment mapping is
inert in the locked rich base medium (nitrogen is supplied by amino-acid
uptakes and ammonium is excreted, so scaling the ammonium bound changes
nothing), which means the nitrogen axis of the environment grid is
effectively unperturbed throughout; we root-caused this during validation
and report it rather than retrofitting the mapping after outcomes.
Second, steady-state flux balance cannot see regulation, damage, or
repair; where the paper's questions touch those processes we say so and
rely on the dynamic and transcriptome-concordance arms rather than on
flux predictions. The paper's claims are exactly the ones that survive
these limitations, and they are stated in Section~\ref{sec:discussion}
with their evidence.

The remainder of the paper is organized as follows.
Section~\ref{sec:related} situates the work in the stress-fermentation
and genome-scale modeling literatures. Section~\ref{sec:methods} gives
the model, the stress layers with their equations and provenance, the
environment space, the analysis stages, the gates, the datasets, and the
pre-registration protocol. Section~\ref{sec:results} reports the staged
results in order. Section~\ref{sec:discussion} interprets the mechanism,
the negatives, and the method boundary, and states the external judge
verdicts and what they changed. Section~\ref{sec:conclusion} concludes.
Appendices carry the locked pre-registration record, parameter tables,
solver-timeout registers, run logs, and the verbatim judge archive.

